\formulary{In Vigilia Circumcisionis Domini et in festo eius}
\label{form:circumcisio}

\rubrica{Ad completorium omnia ut in Nativitate\quickref{form:nativitas}
  excepta antiphona ad Nunc dimittis et collecta.}

\rubrica{Ad Nunc dimittis antiphona}
Glorificamus te Dei Genitrix Virgo:\footnote{\cite[109v]{bp1502}: -Virgo}
quia ex te natus est Christus:
salva omnes qui te glorificant.\footnote{%
  \cite[109v]{bp1502}:
  \rubrica{quae durat usque ad Epiphaniam Domini.}
}\quickref{can:circumcisio_glorificamus}
\edissue{Find a more canonical source.
  The only notated Czech source containing it I could find so far
  is P VI/1, 211v - check it.
  XV A 10 64r knows the antiphon, but references it just
  with an incipit.}

\rubrica{Oratio.}
Deus qui salutis aeternae
beatae Mariae virginitate foecunda
humano generi praemia praestitisti.
tribue quaesumus:
ut ipsam pro nobis intercedere sentiamus:
per quam meruimus auctorem vitae suscipere.
Dominum nostrum Iesum.

\rubrica{Haec sola. Hoc completorium non mutatur usque
  Epiphaniam Domini: excepta sola collecta. videlicet
  \incipit{Deus qui salutis.} Loco cuius dicitur collecta
  \incipit{Praesta quaesumus omnipotens Deus:
    ut natus Salvator}\quickref{nativitas:oratio}
  ut in Nativitate \editorial{Domini}.}

\headingCantus

\cantus{circumcisio_glorificamus}

\edissue{Find a better source}
